\documentclass[aps,pra,reprint,groupedaddress,floatfix]{revtex4-2}
\usepackage{graphicx}
\usepackage{amsmath,amssymb}
\usepackage{placeins}
% Float tuning (reprint two-column carries two large figure*s: D04 2x2, D06 trio).
% Legibility-first: keep panel sizes, relax placement fractions instead.
\renewcommand{\topfraction}{0.9}
\renewcommand{\bottomfraction}{0.8}
\renewcommand{\textfraction}{0.07}
\renewcommand{\floatpagefraction}{0.6}

\graphicspath{{figures/}}

\begin{document}

\title{Variational Quantum State Preparation of the Extended SSH Spin Model\\
on a Digital Quantum Processor}

\author{Xue Qiu}\thanks{These authors contributed equally.}
\author{Yijie Han}\thanks{These authors contributed equally.}
\author{Yijiang Liang}
\author{Jinlong Yu}\thanks{Corresponding author.}
\affiliation{School of Physics and Optoelectronic Engineering, Hainan University, Haikou 570228, China}

\date{\today}

\begin{abstract}
Certifying topological phases on hardware is hard: theorists'
preferred markers are rarely directly measurable, so variational
eigensolvers must build only on directly measurable quantities.
In this work, we prepare the competing trivial, topological, and
antiferromagnetic phases of the dimerized XXZ chain: the three initial
states are the odd-bond singlet product, the topological state of bulk
plus end singlets, and the antiferromagnetic GHZ state, dressed with
reflection-symmetric orbit ansatzes and optimized independently with
the minimum reported. The finite-size phase diagram at $L=8$ is
mapped with the normalized reflection marker $\tilde{Z}_{\mathcal R}$,
resolving the three phases, with gap scaling
confirming controlled finite-size effects. State preparation uses only
directly measurable markers: the AFM structure factor $S(\pi)$ and the
string operator, with premium-qubit screening and tensor-product
readout mitigation bringing the protocol to superconducting hardware.
Hardware and simulator trends
are compared on equal footing via internal $0$--$1$ normalization,
across two dimerizations and both directly measurable markers.
\end{abstract}

\maketitle

\section{Introduction}
\label{sec:intro}

Characterizing topological phases on quantum simulators requires observables that distinguish them from both trivial phases and phases with spontaneous symmetry breaking.
In one dimension, symmetric gapped spin phases can be distinguished by the projective action of symmetry on entanglement~\cite{chen2011classification,schuch2011classifying}.
This SPT distinction differs from intrinsic topological order, for which topological entanglement entropy is a diagnostic~\cite{kitaev2006topological,levin2006detecting}.
The latter has been measured in toric-code states on a superconducting processor~\cite{satzinger2021realizing}.
Here we study the extended Su--Schrieffer--Heeger (SSH) spin model, an interacting dimerized XXZ chain with competing trivial, Haldane-type, and antiferromagnetic (AFM) regimes.

The partial-reflection invariant provides a symmetry-based reference for distinguishing the three regimes in this model~\cite{elben2020manybody}.
Its direct experimental determination requires access to a reflection overlap together with subsystem purities.
Randomized-measurement protocols offer a route to these quantities through mirror-correlated local rotations over multiple measurement settings~\cite{elben2020manybody}.
Direct experimental measurement of the invariant will be addressed in a separate work.

In this work, we instead construct a combined correlation diagnostic $Q$ from two observables accessible in the same $Z$ basis.
The staggered structure factor identifies AFM correlations, while a cell-string correlation provides complementary information on the trivial and topological dimerized regimes.
The interpretation of the string correlation follows from the inserted symmetry and the transformation properties of its endpoints~\cite{pollmann2012detection,haegeman2012order}.
Three reference states fix the affine combination defining $Q$.
We compare its finite-size phase structure with the partial-reflection marker numerically and use it to probe the prepared states experimentally.

Experiments on Rydberg atoms realized a bosonic SSH SPT phase~\cite{deleseleuc2019observation}.
Digital experiments have prepared SPT phases with next-nearest-neighbor spin exchange~\cite{tan2023realizing} and simulated larger spin chains using cloud superconducting devices~\cite{yu2023simulating}.
For variational preparation, Sun \textit{et al.} combined structured initial states and variational layers to study correlated topological phases~\cite{sun2023efficient}.
These experiments establish the setting for exploring which features of the phase diagram can be accessed through a single measurement basis.

To prepare states across the three regimes, we use phase-matched initial states and a reflection-symmetric variational circuit.
Variational quantum eigensolvers combine parameterized circuits with classical optimization to approximate low-energy states~\cite{peruzzo2014variational}.
Superconducting experiments established hardware-efficient implementations for molecules and quantum magnets~\cite{kandala2017hardware}; reviews discuss circuit design, measurement cost, and optimization~\cite{cerezo2021variational,tilly2022variational}.
We optimize the circuits using ideal classical simulations, execute the optimized state-preparation circuits on the Baihua processor, and measure the correlation diagnostics.
The numerical phase map is evaluated at $L=8$; a separate symmetry-sector gap analysis examines finite-size trends along $\delta=-2$.
Readout mitigation follows the assignment-matrix framework reviewed in Refs.~\cite{endo2021hybrid,cai2023quantum}.
The experimental analysis compares measured and simulated $Q$ at $30$ scan points.

The following sections develop the model's symmetry classification, finite-size phase diagram, measured diagnostics, circuit construction, readout mitigation, and comparison with experiment.
The Appendix presents the finite-size gap analysis and qubit-screening results.

\FloatBarrier
\section{Model, symmetries, and topological classification}
\label{sec:model}

\subsection{Dimerized XXZ chain}

We study an open chain of $L$ spin-$1/2$ sites, with $X_j,Y_j,Z_j$ denoting Pauli operators and lattice spacing set to one:
\begin{equation}
H(s,\delta)=\sum_{j=1}^{L-1}t_j
\left[e^{-\delta}(X_jX_{j+1}+Y_jY_{j+1})
+e^\delta Z_jZ_{j+1}\right],
\label{eq:H}
\end{equation}
where $t_{2r-1}=1-s$, $t_{2r}=s$, and $0\leq s\leq1$.
The alternating bonds inherit the pattern of the SSH model~\cite{su1979solitons}.
The free-fermion band classification~\cite{ryu2010topological} provides the noninteracting reference.
Here the spin degrees of freedom and longitudinal interaction lead to a many-body symmetry classification.

To compare conventions with the extended bosonic SSH model of Ref.~\cite{elben2020manybody}, its odd/even couplings are $J=2(1-s)e^{-\delta}$ and $J'=2se^{-\delta}$, and its anisotropy is
\begin{equation}
\Delta=\frac{J_z}{J_{xx}}=e^{2\delta}.
\label{eq:anisotropy}
\end{equation}
Thus $\delta=0$ is the isotropic point, $\delta<0$ is easy plane, and $\delta>0$ is easy axis. This correspondence also connects our model to the interacting bosonic SSH chain realized with Rydberg atoms~\cite{deleseleuc2019observation}. Unless stated otherwise, $L=8$.

\subsection{Symmetries and the variational sector}

The Hamiltonian commutes with the total magnetization $Z_{\rm tot}=\sum_jZ_j$, global spin flip $G_x=\prod_jX_j$, and bond-centered reflection $R:j\mapsto L+1-j$.
It also respects spin time reversal. Continuous spin rotation is $U(1)$ for $\delta\neq0$ and enlarges to $SU(2)$ at $\delta=0$.
For the sizes used here, $L$ is a multiple of four and all three reference states have $Z_{\rm tot}=0$, $G_x=+1$, and $R=+1$.
The conserved operator
\begin{equation}
P=Z_{\rm tot}^2-G_x-R
\end{equation}
packages these conditions: $P=-2$ is equivalent to their simultaneous satisfaction. The ideal variational circuit preserves this sector.

\subsection{Why the two dimerizations are distinct}

A symmetry-protected topological (SPT) phase is defined by equivalence under symmetry-preserving local transformations of gapped states~\cite{gu2009tensor,chen2010local}.
For the onsite symmetries of this chain, the natural unit cell is $(2r-1,2r)$.
Its total spin is $T_r^a=(A_{2r-1}+A_{2r})/2$, with $A=X,Y,Z$ for $a=x,y,z$.
The cell Hilbert space contains both spin-zero and spin-one sectors.
The cell rotations $u_x=e^{i\pi T_r^x}$ and $u_z=e^{i\pi T_r^z}$ commute and generate $D_2\simeq\mathbb Z_2\times\mathbb Z_2$.

For a symmetric, injective matrix product state (MPS), the symmetry acts on its virtual bonds as
\begin{equation}
\sum_n(u_g)_{mn}A^n=e^{i\phi_g}V_g^\dagger A^mV_g,
\qquad V_gV_h=\omega(g,h)V_{gh}.
\label{eq:mps-symmetry}
\end{equation}
The inequivalent projective representations, labeled by $H^2(G,U(1))$ for unitary onsite $G$, distinguish the symmetric gapped phases~\cite{chen2011classification,schuch2011classifying}.
The MPS and Schmidt-decomposition framework is reviewed in Refs.~\cite{schollwock2011density,cirac2021matrix}.
For $D_2$, the gauge-invariant relation
\begin{equation}
V_xV_z=\eta V_zV_x,\qquad \eta=\pm1
\label{eq:projective-index}
\end{equation}
separates the trivial and Haldane-type classes~\cite{pollmann2012detection}.

At $s=0$, singlets lie within cells, giving a product of cell states and $\eta=+1$.
At $s=1$, bulk singlets cross cell boundaries, leaving spin-$1/2$ edge degrees of freedom with $\eta=-1$.
Our topological reference closes the two edge spins into a singlet; this selects one state of the open-chain edge manifold.
This dimer limit illustrates symmetry fractionalization~\cite{verresen2017onedimensional}.
Bond inversion or time reversal can independently protect the Haldane-type distinction~\cite{pollmann2012symmetry}.
The entanglement spectrum, introduced in Ref.~\cite{li2008entanglement}, has symmetry-protected degeneracies in this one-dimensional setting~\cite{pollmann2010entanglement}.
These classification statements concern gapped bulk phases without spontaneous breaking of the protecting symmetry. The AFM region instead has broken spin-flip symmetry in the thermodynamic limit; a finite-chain symmetric cat state can retain the symmetry while displaying large staggered correlations.

\subsection{Partial-reflection diagnostic}

We adopt the bosonic partial-reflection diagnostic of Ref.~\cite{elben2020manybody}.
Partial spatial transformations also enter constructions of many-body invariants for fermionic SPT phases~\cite{shiozaki2017manybody}.
For a central interval $I=I_1\cup I_2$, let $\rho_I$ be its reduced density matrix and $\mathcal R_I$ reverse its sites. We evaluate
\begin{equation}
\tilde Z_{\mathcal R}
=\frac{\operatorname{Tr}(\rho_I\mathcal R_I)}
{\sqrt{\left[\operatorname{Tr}(\rho_{I_1}^2)
+\operatorname{Tr}(\rho_{I_2}^2)\right]/2}}.
\label{eq:reflection-marker}
\end{equation}
For $L=8$, $I=(3,4,5,6)$ and its halves are $(3,4)$ and $(5,6)$.
In the symmetric gapped regimes, the quantized limits $+1$ and $-1$ require intervals large compared with the correlation length; the AFM limit is zero in the corresponding model~\cite{elben2020manybody}.
We use the four-site interval to evaluate the finite-size phase map.

We compute Eq.~\eqref{eq:reflection-marker} from numerical wave functions.
The expression combines off-diagonal matrix elements of $\mathcal R_I$ with subsystem purities.
Randomized measurements with suitable basis rotations provide an experimental route to these quantities~\cite{elben2022toolbox}.

\FloatBarrier
\section{Finite-size phase diagram}
\label{sec:phasediagram}

Figure~\ref{fig:D04}(d) shows the normalized partial-reflection marker over the $(s,\delta)$ plane for $L=8$. Each grid point is computed independently from the corresponding ground state. At degenerate points, we retain the lowest-energy eigenvector returned by the solver and record a degeneracy flag in the dataset. The map defines finite-size parameter regions, representative points, and reference parameters for variational preparation. The Appendix examines the size dependence of the gap along $\delta=-2$. The correlators in panels (a)--(c) are discussed in Sec.~\ref{sec:heatmaps}.


\begin{figure}[htbp]
\centering
\includegraphics[width=\columnwidth]{{../../figs/exp02/exp02_D04}}
\caption{Topological heatmaps over the $(s,\delta)$ plane:
(a) $S(\pi)$, (b) string operator,
(c) $Q$, (d) $\tilde{Z}_{\mathcal R}$. \label{fig:D04}}
\end{figure}

\FloatBarrier
\section{Phase markers}
\label{sec:markers}

\subsection{Structure factor and the AFM label}
\label{sec:sq}

To identify the antiferromagnetic phase we use the structure factor
\begin{equation}
S(q)=\frac{1}{L}\sum_{i,j}e^{iq(i-j)}\langle Z_iZ_j\rangle,
\end{equation}
where $\langle Z_iZ_j\rangle$ is the ground-state correlator and $L=8$
throughout this section.
The structure factor is the Fourier transform of spin correlations; $S(\pi)$ quantifies staggered correlations in the finite chain. The $S(q)$ curves at one representative point in each
of the three phases (labels trivial / topo / AFM, coordinates in the
dataset) on the shared $q$ grid (including $q=\pi$) are shown in
Fig.~\ref{fig:D03}.
The $q=\pi$ component is used as the AFM diagnostic throughout the
hardware-facing protocol. It is evaluated from $Z_iZ_j$ correlations obtained by computational-basis measurements. Hardware $p=1$ measurements at the three points (Fig.~\ref{fig:D03}(b)) track the simulated curves in Fig.~\ref{fig:D03}(a), with the AFM $\pi$ peak clearly resolved despite noise suppression.

\begin{figure}[htbp]
\centering
\includegraphics[width=\columnwidth]{{../../figs/exp02/exp02_D03}}
\caption{Structure factor $S(q)$ at three phase-representative points (one
per phase) on the shared $q$ grid: (a) simulator, (b) hardware $p=1$
measurements at the same points ($1024$ shots with readout mitigation),
tracking the simulated curves; the values at $q=\pi$ compare staggered
correlations at the representative points.
\label{fig:D03}}
\end{figure}

\FloatBarrier
\subsection{Correlation heatmaps}
\label{sec:heatmaps}

Panels (a)--(c) of Fig.~\ref{fig:D04} show $S(\pi)$, the string correlation $O_{\rm str}$, and the combination
\begin{equation}
Q=\frac{4}{3}+2O_{\rm str}-\frac{1}{6}S(\pi).
\end{equation}
All three quantities are evaluated on the same parameter grid as the reflection marker, allowing comparison of the parameter dependence of staggered correlations, nonlocal string correlations, and their combination.

The general form of the string correlator is
\begin{equation}
O^z_{\rm str}(i,j)=\Big\langle S^z_i\exp\!\Big(i\pi\!\sum_{k=i+1}^{j-1}S^z_k\Big)S^z_j\Big\rangle,
\end{equation}
where $S^z_k=Z_k/2$. The spin-dependent phase factor probes hidden order within the interval, providing a diagnostic of valence-bond phases~\cite{dennijs1989preroughening,perezgarcia2008string}. For a finite chain, the measured quantity is the correlation between selected endpoints. On our $L=8$ chain we use $d=3$ with summed endpoint pairs, $O_{\rm str}=\langle(Z_1+Z_2)(Z_3Z_4)(Z_5Z_6)(Z_7+Z_8)\rangle$ (the two minus signs cancel for $L=8$); the heatmaps and the hardware protocol use $\hat S(\pi)=S(\pi)/8$ and $\hat O_{\rm str}=-O_{\rm str}$.

Both $S(\pi)$ and the $z$-string correlation can be evaluated from computational-basis bitstring statistics and serve as the main observables in our hardware protocol. The partial-reflection marker is used for the numerical phase diagram; measurement schemes for many-body invariants and randomized measurements are discussed in Refs.~\cite{elben2020manybody,elben2022toolbox}. The observable scales used in the hardware--simulator comparison are specified in Sec.~\ref{sec:comparison}.

\FloatBarrier
\section{Variational state preparation}
\label{sec:ansatz}

\subsection{Optimization targets}
\label{sec:points}

Variational optimization targets each point of the two scan lines, $\delta=-2$ and $\delta=0.85$, on the $99$-point $s$ grid ($s=0.01$--$0.99$): at every grid point the three initial states are optimized independently and the lowest of the three minima is reported, so the target Hamiltonians sample the trivial, topological, and AFM regions mapped in Fig.~\ref{fig:D04}(d). The lowest-energy initial-state branches along the two lines are shown by the background bands in Fig.~\ref{fig:D08D09}(a,b).

\subsection{Variational circuit structure}

Our variational circuits follow the variational eigenvalue solver
paradigm~\cite{peruzzo2014variational}: prepare an initial state, then
apply a parameterized ansatz. The classical loop minimizes
\begin{equation}
E(\theta)=\langle\psi(\theta)|H|\psi(\theta)\rangle,
\end{equation}
where $|\psi(\theta)\rangle$ is the circuit output at parameters
$\theta$. At each target point, the three reference states are optimized independently at every depth $p=1,2,3$.
We screen random parameter vectors together with the zero vector, refine the low-energy candidates with local L-BFGS, and retain the lowest-energy converged result in each branch.
The minimum across the three branches determines $E^*$ and its parameters.
The symmetry-preserving circuit retains the initial symmetry quantum numbers.
The scan runs point by point along the $\delta=-2$ and $\delta=0.85$ lines. Variational algorithms of this form are surveyed in
Refs.~\cite{cerezo2021variational,tilly2022variational}, which collect
cost-function design, ansatz families, and best practices for VQE.

The ideal-state optimization is performed classically; the processor executes the resulting circuits to measure correlators.
For a normalized trial state in a sector with a unique ground state, spectral decomposition gives
\begin{equation}
1-|\langle\psi_0|\psi(\theta)\rangle|^2
\leq \frac{E(\theta)-E_0}{\Delta_{\rm sec}}.
\label{eq:energy-fidelity}
\end{equation}
Equation~\eqref{eq:energy-fidelity} relates the trial-state overlap to the energy residual through the sector gap $\Delta_{\rm sec}$.

Our bond rotations form a Hamiltonian variational ansatz: the generators follow the interaction terms and their angles are optimized independently~\cite{wecker2015progress}.
The XXZ setting has been studied within this family~\cite{wiersema2020exploring}.
Building symmetry into state-preparation circuits restricts the search to relevant sectors~\cite{gard2020efficient}.
Different reference states supply different initial entanglement patterns, a relevant resource in variational preparation of nontrivial states~\cite{ho2019efficient}.
The three reference states share the same finite-chain quantum numbers and provide distinct entanglement structures; their optimized energies determine the branch selected at each parameter point.

Trainability also depends on circuit and cost structure.
Random circuits approaching unitary designs can exhibit exponentially small gradients~\cite{mcclean2018barren}, and local versus global costs can behave differently even at shallow depth~\cite{cerezo2021cost}.
Identity-block initialization~\cite{grant2019initialization} and model-based circuit design~\cite{wiersema2020exploring} motivate structured starting points.
Depth-dependent optimization transitions~\cite{campos2021abrupt} and noise-induced gradient suppression under specified noise models~\cite{wang2021noise} further characterize the roles of circuit depth and noise in variational optimization.

Let $M=L/2$. The three reference initial states lie in the $P=-2$ sector and are shown in Fig.~\ref{fig:D06}. They are paired with an
orbit ansatz in which mirror-symmetric bond pairs share rotation angles
(odd bonds $O_j\leftrightarrow O_{M+1-j}$, even bonds
$E_j\leftrightarrow E_{M-j}$). The ansatz is built in two steps:
first fix the entanglement structure with a phase-matched initial state,
then dress all bonds with rotations tied by the reflection symmetry, so
the parameter count grows with the number of bond orbits rather than
bonds. Each sublayer evolves under ($\alpha$ labels the orbit,
$\langle kl\rangle$ a bond):
\begin{equation}
U_{o}=\prod_{\alpha}\prod_{\langle kl\rangle\in\mathcal{O}_{\alpha}}e^{-i\theta_{o,\alpha}(X_{k}X_{l}+Y_{k}Y_{l})/2}e^{-i\varphi_{o,\alpha}Z_{k}Z_{l}/2},
\end{equation}
\begin{equation}
U_{e}=\prod_{\alpha}\prod_{\langle kl\rangle\in\mathcal{E}_{\alpha}}e^{-i\theta_{e,\alpha}(X_{k}X_{l}+Y_{k}Y_{l})/2}e^{-i\varphi_{e,\alpha}Z_{k}Z_{l}/2}.
\end{equation}
On the same bond $XX{+}YY$ commutes with $ZZ$, and disjoint bonds within
a sublayer commute; the odd and even sublayers do not commute with each
other, so both are kept. Each bond is implemented as an
$R_{XX}(\theta)R_{YY}(\theta)R_{ZZ}(\varphi)$ gate sequence. Each orbit
carries one $(\theta,\varphi)$ pair, with $L$ parameters per layer for the $L=8$ chain considered here (halved
at $\delta=0$ where $\theta=\varphi$); the ansatz preserves
$Z_{\rm tot}$, $\prod X$, and $R$, staying in the $P=-2$ sector. The
sublayer overlapping the initial-state pairing acts last: even-first
for trivial and AFM initial states, odd-first for topological. One schematic per phase is shown in Fig.~\ref{fig:D06}(a)--(c): initial state
plus required ansatz circuit for the trivial, topological, and
AFM cases.
For ideal circuits, the trial states remain in the target symmetry sector. Section~\ref{sec:comparison} compares measured and simulated correlators.
Figure~\ref{fig:D08D09}(a,b) shows the $p=1$ cost functions
against the exact ground-state energy.

\begin{figure*}[htbp]
\centering
\includegraphics[width=\textwidth]{{../../figs/exp04/exp04_D06}}
\caption{Initial states and required ansatz circuits for the three phases:
(a) trivial; (b) topological; (c) AFM. The dashed box $U^{(k)}$ is the $k$th variational layer;
$\theta^{(k)}_j$ ($\varphi^{(k)}_j$) is the $XX{+}YY$ ($ZZ$) angle shared by
the $j$th mirror-symmetric bond orbit in layer $k$ (odd
$O_j\leftrightarrow O_{M+1-j}$, even $E_j\leftrightarrow E_{M-j}$;
$\theta=\varphi$ at $\delta=0$). Even-bond sublayer first in (a) and (c),
odd-bond sublayer first in (b). \label{fig:D06}}
\end{figure*}

\begin{figure*}[htbp]
\centering
\includegraphics[width=\textwidth]{{../../figs/exp06/exp06_D08D09}}
\caption{Cost functions at $p=1$ and the experiment--simulation comparison.
(a),(b) Cost functions at $p=1$ versus $s$ for (a) $\delta=-2$ and (b) $\delta=0.85$:
$E_0$ (exact, blue solid) with the three initial-state costs
(triv green dashed, topo orange dotted, AFM black dashdot);
background bands identify the initial-state branch minimizing the $p=1$ cost at each grid point; circles identify the branch with the lowest mean optimized cost within each interval.
(c),(d) Combined quantity $Q$ for (c) $\delta=-2$ and (d) $\delta=0.85$:
simulation $p=1$ (black) and $p=2$ (gray) fixed-scale curves with
experimental window scatter ($p=1$, 16-unit window of the 500-unit round; colors indicate
the lowest-energy initial-state branches in (a),(b)).
\label{fig:D08D09}}
\end{figure*}

\FloatBarrier
\section{Readout error mitigation}
\label{sec:readout}

We model classical readout misclassification by an assignment matrix,
\begin{equation}
\mathbf p_{\rm meas}=\mathsf A\mathbf p_{\rm ideal},
\qquad \mathsf A_{ab}=P(a|b).
\label{eq:assignment}
\end{equation}
Here $\mathbf p_{\rm ideal}$ is the computational-basis distribution of the prepared state immediately before readout~\cite{cai2023quantum}.
Its columns are estimated by preparing computational-basis states.
For qubit $i$, let $\epsilon_{0,i}=P(1|0)$ and $\epsilon_{1,i}=P(0|1)$.
The independent-error approximation is
\begin{equation}
\mathsf A=\bigotimes_{i=1}^L
\begin{pmatrix}
1-\epsilon_{0,i}&\epsilon_{1,i}\\
\epsilon_{0,i}&1-\epsilon_{1,i}
\end{pmatrix}.
\label{eq:tensor-readout}
\end{equation}
This tensor-product model contains $2L$ single-qubit misclassification probabilities~\cite{bravyi2021mitigating}.
One pair of all-zero and all-one preparation circuits calibrates the single-qubit marginals in each acquisition unit.
The platform's built-in readout correction is disabled.

We first normalize the counts by their actual total to obtain $\mathbf p_{\rm meas}$.
When the calibrated matrix is invertible, linear inversion gives $\mathbf v=\mathsf A^{-1}\mathbf p_{\rm meas}$.
We obtain a probability distribution by projecting the joint inverse vector onto the probability simplex in Euclidean distance~\cite{maciejewski2020mitigation}:
\begin{equation}
\mathbf p_{\rm mit}
=\underset{\mathbf p\ge0,\;\mathbf 1^{\mathsf T}\mathbf p=1}
{\operatorname{argmin}}\lVert\mathbf p-\mathbf v\rVert_2^2.
\label{eq:readout-projection}
\end{equation}
The solution is $p_{{\rm mit},b}=\max(v_b-\tau,0)$, where a common threshold $\tau$ enforces unit sum.
We determine this threshold by sorting, then compute the correlations from the projected distribution.
The projection is nonlinear and can introduce finite-sampling bias.

We recalibrate in each acquisition unit to track temporal variation.
Joint acquisition uses one calibration pair on the selected chain mapping to construct a tensor-product readout model shared by the branches.
The separate measurements at the three representative points each use calibration on their own physical mapping.
The hardware analysis reports the resulting correlations and their inter-unit standard errors.

\FloatBarrier
\section{Hardware--simulator comparison}
\label{sec:comparison}

Figure~\ref{fig:D08D09}(c,d) compares the combined quantity $Q$ from simulations at $p=1,2$ and hardware at $p=1$. Measurements on the Baihua processor are organized into repeated acquisition units. Each unit contains $15$ values of $s$ along two scans, $\delta\in\{0.85,-2\}$, giving $30$ parameter points with $1024$ shots per point and one additional pair of readout-calibration circuits. A single round of $500$ units is analyzed. Error bars on hardware points indicate standard errors across acquisition units.

The observables $\hat S(\pi)$ and $\hat O_{\rm str}$ are evaluated from the same readout-mitigated distributions. Hardware hats keep the same sign convention ($\hat O=-O_{\rm str}$) and are divided by the selected window's shared maxima: the top-1 maxima of the window-pooled $30$-point means, refitted for each candidate window. For the selected window these are $S_{\max}=4.51$ and $O_{\max}=0.090$. Simulations use fixed scales, $\hat S=S/8$ and $\hat O=-O_{\rm str}$, and the $p=1$ simulation is the reference curve for window selection ($p=2$ is shown for comparison only). Both use
\begin{equation}
Q=\frac{4}{3}-2\hat O-\frac{4}{3}\hat S.
\end{equation}
With the fixed simulation scales, this expression agrees with the definition in Sec.~\ref{sec:heatmaps}. The combined $Q$ values are used directly, with no post-combination renormalization, to evaluate the distance metric $d_{\rm raw}$, the mean absolute $Q$ distance over the $30$ points, and the Pearson correlation coefficient $r$.

Among the compared even-length windows, the best window was $[226,242)$, comprising $16$ acquisition units, with $d_{\rm raw}=0.30$ and $r=0.85$. Candidate windows are all contiguous even-length subsegments (length at least $10$) of the $500$-unit sequence, each with its own refitted endpoints; the best window maximizes $r$ against the $p=1$ simulation, with ties broken by smaller $d_{\rm raw}$, then longer length, then earlier start. These metrics are report-only descriptions of curve agreement under the selected window and its normalization.

Along $\delta=-2$, hardware results exhibit a parameter dependence similar to the simulated $p=2$ curve throughout the scan. At $\delta=0.85$, the first half follows a similar trend, while in the topological region hardware $Q$ turns negative and reaches lower values than either simulation as $s\to1$. A resolved string signal occurs at $s=0.99$: pooling all $N=500$ units, the hardware $\hat O_{\rm str}$ mean is approximately $5\%$ of the theoretical value, $9\sigma$ away from zero where $\sigma$ is the inter-unit standard error; the remaining points are consistent with zero within measurement uncertainty.

\FloatBarrier
\section{Conclusion}
\label{sec:conclusion}

We explored the phase diagram of the extended SSH spin model through a combined correlation diagnostic accessible in a single $Z$ measurement basis.
The staggered structure factor identifies AFM correlations, while the cell-string correlation supplies complementary information on the trivial and topological dimerized regimes.
Their reference-state-calibrated combination $Q$ captures a three-region structure similar to the partial-reflection marker in the $L=8$ numerical map.
A two-site-cell symmetry description connects the dimer limits to distinct projective representations, and the sector-gap trend along $\delta=-2$ is consistent with the gapless uniform-chain prediction.

We prepared states with phase-matched initial states and a reflection-symmetric variational circuit, optimized classically and executed on the Baihua processor.
The experimental protocol combines cluster-state-based qubit screening, tensor-product readout correction, and measurements of both constituent correlators from the same computational-basis samples.
For the displayed data, the $30$-point experiment--simulation comparison uses $16\,384$ shots per point and gives $r=0.86$ with fitted marker scales.

The measurement strategy relies on site-resolved spin or equivalent occupation readout.
Optical-lattice imaging has already accessed nonlocal string correlations~\cite{endres2011observation}, and optical-tweezer Rydberg arrays provide site-resolved readout in a bosonic SSH realization~\cite{deleseleuc2019observation}.
These capabilities motivate transferring the correlation diagnostic to such simulators when the corresponding spin encoding and reference-state calibration are available.

\FloatBarrier
\appendix

\section{Finite-size gap analysis}
\label{sec:gap}

We take the gap in the $P=-2$ symmetry sector,
\begin{equation}
\Delta_{\rm sec}(L,s)=E_1^{(P=-2)}(L,s)-E_0^{(P=-2)}(L,s),
\end{equation}
where $E_0^{(P=-2)}$ ($E_1^{(P=-2)}$) is the ground-state (first-excited-state)
energy in that sector.
The sector gap at $\delta=-2$ for $L=8,12,16,20,24$ is shown in
Fig.~\ref{fig:D02a}(a). Panel~(b) zooms to $s\in[0.53,0.65]$
(Fig.~\ref{fig:D02b}), marking the minimum of each curve: as $L$
grows the minima move toward $s=0.5$ and their values keep
decreasing, consistent with convergence to $s=0.5$ in the
thermodynamic limit. The gap at $s=0.5$ versus $1/L$
is shown in Fig.~\ref{fig:D02c}(c); the least-squares quadratic fit in $1/L$ over the five sizes has intercept $a=-0.017$ in units of the exchange coupling, small against the gap scale and consistent with
the gap-closing trend. These results characterize the finite-size gap evolution near $s=0.5$ at $\delta=-2$.
Ground states throughout are obtained by exact diagonalization
(Lanczos), and the minima analysis follows standard finite-size
practice for quantum spin systems, for which we follow the lecture
notes of Ref.~\cite{sandvik2010computational}.

\begin{figure*}[htbp]
\centering
\includegraphics[width=\textwidth]{{../../figs/exp03/exp03_D02}}
\caption{Finite-size gap analysis at $\delta=-2$: (a) symmetry-sector gap
$\Delta_{\rm sec}$ for $L=8,12,16,20,24$; (b) zoom to
$s\in[0.53,0.65]$ with curve minima marked, moving toward $s=0.5$ as
$L$ grows; (c) gap at $s=0.5$ versus $1/L$ for five
sizes with the least-squares quadratic fit in $1/L$; the square
marks the fitted intercept $a=-0.017$.
\label{fig:D02a}\label{fig:D02b}\label{fig:D02c}}
\end{figure*}

\section{Premium-qubit screening}
\label{sec:premium}

To run on hardware we screen candidate chains (rings) with two
quantities: the stabilizer mean, probing two-qubit gate quality, and
the readout fidelity, probing measurement quality.
The ideal output is the cluster state
$|C\rangle=\prod_{\langle ij\rangle}CZ_{ij}|+\rangle^{\otimes n}$,
where $\langle ij\rangle$ runs over the edges of the candidate subgraph;
its stabilizers $S_i$ (two-body $X_0Z_1$, $Z_{n-2}X_{n-1}$ at open ends,
three-body $Z_{i-1}X_iZ_{i+1}$ otherwise) have expectation $+1$ in the
absence of errors, so deviations of $\langle S_i\rangle$ from $1$ characterize combined state-preparation and measurement errors.
Each candidate is characterized by $4$ benchmark circuits: two
stabilizer circuits measuring the even- and odd-index stabilizers
separately (the stabilizers commute, while the two groups use different local measurement bases),
and two readout circuits without $CZ$ gates yielding the readout
fidelity $F_{\rm ro}$ from the mean of $P(00\cdots0)$ and
$P(11\cdots11)$. The score
\begin{equation}
{\rm score}=0.8\,\bar{S}+0.2\,F_{\rm ro},
\end{equation}
where $\bar{S}=n^{-1}\sum_i\langle S_i\rangle$ is the stabilizer mean,
ranks the candidates, weighting two-qubit gate quality over readout.
Cluster states~\cite{raussendorf2001oneway}
are the ideal outputs whose stabilizers we measure.
The $8$-chain and $10$-ring benchmarks differ (the ring includes a
closing $CZ$ and has three-body stabilizers throughout), so candidate rankings are computed separately within each geometry. The benchmark evaluates candidate subgraphs through cluster-state preparation and local measurements on the specified connectivity. The resulting rankings and selected qubits are reported below.

\section{Premium-qubit rankings}
\label{sec:rankings}

The screening on the Baihua processor yields three rankings. In the
$8$-chain geometry the top ten candidates are headed by the
$8$-qubit chain on qubits 77 through 70 with score $0.7801$ (one further
candidate tied within $0.02$). In the $10$-ring geometry the top ten are
headed by the $10$-qubit ring on qubits 69--73 and 82--86 with score
$0.5246$ (two further candidates tied within tolerance). Among the ten
$8$-subchains of the top-ranked ring, the subchain on qubits
69--73 and 82--84 is first, with score $0.7108$ taken from the
chain-benchmark measurements. Subsequent hardware experiments use the
leading candidate of each ranking.




\FloatBarrier
\bibliographystyle{unsrtnat}
\bibliography{references}

\end{document}
